\section{Introduction}

The imminent arrival of 6G networks promises to empower next-generation edge intelligence with ultra-low latency, high reliability, and context-aware services—including XR, autonomous systems, and tactile internet applications~\cite{Lou2022_CostEffectiveScheduling, Wang2025_AverageReliabilityOptimal}. These use cases demand fine-grained orchestration of computing and data services across heterogeneous edge nodes. Traditional cloud-based approaches cannot meet such requirements due to latency bottlenecks, reinforcing the importance of edge computing to bring processing closer to end users~\cite{Xu2023_FedLC}.

Nevertheless, the dynamic and resource-constrained nature of edge environments—combined with energy, privacy, and workload variability—introduces substantial orchestration challenges~\cite{Yuan2025}. Recent studies have investigated reinforcement learning (RL) for adaptive orchestration, allowing real-time learning of placement and migration strategies~\cite{Yuan2025}. However, centralized RL methods often suffer from scalability issues and cannot preserve user privacy. Furthermore, most RL-based orchestration agents disregard semantic content in service tasks, leading to suboptimal placements in latency-sensitive or context-dependent scenarios such as XR or real-time analytics~\cite{Yu2025}.

Federated reinforcement learning (FRL) provides a promising alternative by distributing policy learning across edge nodes without sharing raw data~\cite{Chen2025_FRL}. While FRL frameworks have shown promise in handling privacy and scalability, they frequently overlook two critical dimensions: semantic task awareness and network topology encoding. In addition, many studies perform limited evaluations in synthetic environments with minimal orchestration realism or baseline diversity.

Modern edge workloads frequently include semantic labels such as task intent or content category, which—if properly modeled—can significantly enhance placement fidelity and decision stability~\cite{Wu2024_FRLEdgeComputing}. At the same time, orchestration effectiveness is shaped by the network structure, including link quality, proximity, and cluster connectivity. Yet most FRL methods treat edge nodes as unstructured agents. While graph neural networks (GNNs) offer an effective means to encode such topology~\cite{Sagar2025_HAFedRL}, their integration into federated, real-time orchestration remains limited.

Moreover, frequent global synchronization in FL and FRL systems often results in high communication overhead and latency. Hierarchical FL frameworks reduce this burden by performing local aggregation within clusters before periodic global updates~\cite{Ma2024_HierarchicalFL}. However, few approaches combine hierarchical control with semantic embeddings and GCN-based topology encoding in a unified orchestration framework. To the best of our knowledge, no prior study jointly evaluates these dimensions in a full-stack federated orchestration simulator using realistic simulation traces with detailed instrumentation. Large-scale scenarios (1,000–10,000 nodes) are analyzed through mathematical extrapolation from empirically validated base simulations, allowing principled scalability analysis without incurring prohibitive computational overhead.

\noindent\textbf{Contributions.} We propose \textbf{FedSemGNN}, a federated reinforcement learning framework designed for semantic-aware, topology-informed, and hierarchical orchestration in 6G edge environments. The key contributions of this work include:

\begin{itemize}
	\item \textbf{Semantic task embedding with online learning:} Service requests are embedded using pre-trained Sentence-BERT (\texttt{all-MiniLM-L6-v2}) in 128-dimensional semantic spaces and matched against node capabilities using cosine similarity, enabling intent-aware placement that goes beyond numerical resource matching. Online semantic learning with elastic weight consolidation ($\lambda_{\text{ewc}} = 0.4$) enables real-time adaptation to evolving workload characteristics.
	
	\item \textbf{GCN-based topology encoding:} A two-layer graph convolutional network with symmetric normalization captures structural relationships—such as bandwidth, affinity, and proximity—allowing placement decisions that consider network-wide constraints and multi-hop dependencies across mesh-connected edge infrastructure.
	
	\item \textbf{Hierarchical federated PPO with dual-level synchronization:} Edge servers are organized into 10 semantic-topological clusters that perform localized PPO training with two-level aggregation: intra-cluster synchronization every $K_1 = 10$ epochs and inter-cluster global synchronization every $K_2 = 50$ epochs. This reduces communication overhead while maintaining convergence stability through differential privacy ($\sigma_{\text{DP}} = 1.0$) and gradient clipping (norm 0.5).
	
	\item \textbf{Compact neural architecture for ultra-low latency:} The PPO agent employs a lightweight actor-critic network with only 16 hidden units per layer, enabling sub-millisecond orchestration decisions (0.36 ms per step) while maintaining high reward performance through prioritized experience replay (buffer size 20,000) and entropy regularization ($\beta = 0.01$).
	
	\item \textbf{System-level instrumentation in EdgeSimPy~\cite{edgesimpy}:} We implement FedSemGNN in a realistic edge computing simulator with fine-grained logging of control-plane latency, reward trajectories, semantic fidelity, power consumption, communication footprint, and migration dynamics across 6 heterogeneous edge servers and 500 mobile users over 500 timesteps.
	
	\item \textbf{Comprehensive empirical validation with scalability analysis:} Using a 6-server baseline testbed with 500 orchestration steps, FedSemGNN achieves:
	\begin{itemize}
		\item \textbf{Normalized reward of 0.91}—2.9$\times$ higher than Random (0.31) and 1.9$\times$ higher than GreedyPower (0.47),
		\item \textbf{Ultra-low latency of 0.36 ms per decision}—315$\times$ faster than FlatFedPPO (114.25 ms) and 222$\times$ faster than HierFedPPO (80.71 ms),
		\item \textbf{Communication efficiency of 1.394 reward per MB}—demonstrating superior reward-to-bytes ratio through hierarchical aggregation,
		\item \textbf{100\% semantic fidelity with zero migrations}—compared to 6,840 migrations in FlatFedPPO and 1,181 in HierFedPPO,
		\item \textbf{Energy consumption of 72.1 W}—36.2$\times$ lower than FlatFedPPO (2607.9 W) and 14.7$\times$ lower than HierFedPPO (1057.4 W).
	\end{itemize}
	Scalability to 1,000–10,000 nodes is evaluated through validated mathematical extrapolation using power-law and logarithmic scaling models fitted to empirical base results.
\end{itemize}

Unlike previous methods that treat semantics or topology heuristically, FedSemGNN unifies both within a scalable, federated RL loop with mathematically grounded reward optimization. The reward function $R_t = \max(0, 10{,}000 - (P_t + 0.01 \cdot L_t))$ jointly minimizes power consumption and latency while implicitly encouraging semantic fidelity through state representation design. By embedding semantic structure into the learning process and encoding topological constraints through GNNs, FedSemGNN supports responsive orchestration with minimal latency and communication cost. The system demonstrates real-time operation under 0.5 ms per step, offering a viable foundation for control-plane design in 6G edge computing.

\noindent\textbf{Paper organization.} The remainder of this paper is structured as follows: Section~\ref{sec:related} reviews related work in federated RL, semantic scheduling, and GNN-based orchestration. Section~\ref{sec:system_architecture} presents the FedSemGNN system architecture including semantic embedding, GNN topology encoding, and hierarchical PPO coordination. Section~\ref{sec:methodology} details the proposed methodology with mathematical formulations for reward optimization, policy learning, and federated aggregation. Section~\ref{sec:results} discusses empirical evaluations across baseline comparisons, ablation studies, and scalability analysis. Section~\ref{sec:limitations_future} outlines limitations and future research directions. Section~\ref{sec:conclusion} concludes the paper.
